紧算子的另一个有用性质是, 它将弱收敛序列映为强收敛序列

\begin{Theorem}
\label{thm:compact-operator-weak-to-strong}
假设 $A\in\mathcal{L}(X,Y)$ 为紧算子, 且
\begin{equation}\label{eq:compact-operator-weak-convergence-hypothesis}
x_n\rightharpoonup x\text{ (weakly) in }X.
\end{equation}
那么
\begin{equation}\label{eq:compact-operator-strong-image-convergence}
A(x_n)\to A(x)\text{ (strongly) in }Y.
\end{equation}

\end{Theorem}

\begingroup
\postdisplaypenalty=10000
\renewcommand{\endProof}{\par\nobreak\hfill$\blacksquare$\par\vspace{0.1cm}}
\begin{Proof}
第一步是证明
\begin{equation}\label{eq:bounded-operator-weak-image-convergence}
A(x_n)\rightharpoonup A(x)\text{ (weakly) in }Y.
\end{equation}
给定 $f\in Y^*$. 需要证明
\begin{equation}\label{eq:bounded-operator-weak-image-functional-limit}
\lim_{n\to\infty}f(A(x_n))=f(A(x)).
\end{equation}
为此, 用下式定义 $g:X\to\mathbb{R}$:
\begin{equation}\label{eq:bounded-operator-composed-dual-functional}
g(z)=f(A(z)),\quad z\in X.
\end{equation}
由于 $f$ 和 $A$ 都是线性的, $g$ 为线性映射, 且 $g$ 有界, 因为
\begin{equation}\label{eq:bounded-operator-composed-functional-norm-estimate}
|g(z)|=|f(A(z))|\leq\|f\|\|A(z)\|\leq\|f\|\|A\|\|z\|.
\end{equation}
因而 $g\in X^*$. 又由于 $x_n\rightharpoonup x$ 于 $X$ 中成立, 有\[\begin{aligned}
\lim_{n\to\infty}f(A(x_n))&=\lim_{n\to\infty}g(x_n)\\
&=g(x)\\
&=f(A(x)).
\end{aligned}\]由于 $f$ 任意, $A(x_n)\rightharpoonup A(x)$.

现在假设 $A(x_n)$ 在 $Y$ 中不强收敛到 $A(x)$. 那么存在 $\varepsilon>0$ 及子序列 $A(x_{n_k})$, 使得
\begin{equation}\label{eq:compact-operator-strong-convergence-contradiction-sequence}
\|A(x_{n_k})-A(x)\|\geq\varepsilon.
\end{equation}
现在, 由于 $x_n$ 弱收敛到 $x$, $x_{n_k}$ 也如此. 由于 $x_{n_k}$ 弱收敛, 它有界. 因而由 $A$ 的紧性, $A(x_{n_k})$ 有强收敛子序列. 然而, 由于强收敛蕴含弱收敛, 且弱极限唯一, 这一子序列必收敛到 $A(x)$. 但这与 \eqref{eq:compact-operator-strong-convergence-contradiction-sequence} 矛盾, 证明完毕.
\end{Proof}
\endgroup

将这一结果与 Theorems~\ref{thm:sobolev-continuous-target-compact-embedding} 和 \ref{thm:compact-successive-sobolev-orders} 结合, 得到以下推论.

\begin{Corollary}
\label{cor:compact-sobolev-weak-to-uniform}
假设 $k>m/2$, $\Omega\subset\mathbb{R}^m$ 有界且具有 $k$ 阶延拓性质. 设
\begin{equation}\label{eq:compact-sobolev-weak-to-uniform-hypothesis}
u_n\rightharpoonup u\text{ (weakly) in }H^k(\Omega).
\end{equation}
那么
\begin{equation}\label{eq:compact-sobolev-weak-to-uniform-conclusion}
u_n\to u\text{ uniformly on }\overline\Omega.
\end{equation}

\end{Corollary}

\begin{Corollary}
\label{cor:compact-sobolev-weak-to-lower-strong}
假设 $k$ 为非负整数, $\Omega\subset\mathbb{R}^m$ 有界且具有 $k+1$ 阶延拓性质. 设
\begin{equation}\label{eq:compact-sobolev-successive-order-weak-hypothesis}
u_n\rightharpoonup u\text{ (weakly) in }H^{k+1}(\Omega).
\end{equation}
那么
\begin{equation}\label{eq:compact-sobolev-successive-order-strong-conclusion}
u_n\to u\text{ (strongly) in }H^k(\Omega).
\end{equation}

\end{Corollary}
